\section{Marco conceptual}

\begin{itemize}

\item \textbf{Aprender a aprender:} En el 2006, el parlamento europeo y el Consejo de Educación definieron aprender a aprender como la capacidad de perseguir y persistir en el aprendizaje, organizar el propio aprendizaje, incluyendo la gestión efectiva del tiempo y la información, tanto de manera individual como en grupo. Esta competencia abarca el reconocimiento del propio proceso de aprendizaje y sus necesidades, identificar las oportunidades disponibles y la capacidad de superar obstáculos para aprender con éxito. Implica adquirir, procesar y asimilar nuevos conocimientos y habilidades, así como buscar y utilizar orientación \cite{GarciaVazquez2015}.
\\
\\
Aprender a aprender motiva a los individuos a basarse en aprendizajes previos y experiencias de vida para aplicar conocimientos y habilidades en diversos contextos, ya sea en el hogar, en el trabajo, en la educación o la formación. La motivación y la confianza son elementos esenciales para esta competencia.

\item \textbf{Aprendizaje activo:} Se fundamenta en la teoría del Constructivismo, la cual destaca que los alumnos construyen su propio conocimiento. Este enfoque de enseñanza implica que los estudiantes participen activamente en el proceso de aprendizaje, desarrollando su conocimiento y comprensión a través de oportunidades de aprendizaje diseñadas por sus docentes \cite{cambridge2024}.

\item \textbf{Escucha activa:} El término “escucha activa” fue introducido en 1957 por Carl Rogers y Richard Farson. También puede ser referido como escucha reflexiva. Es una habilidad de comunicación que implica no limitarse a solo escuchar las palabras que otra persona dice, si no escuchar con la intención de entender completamente a la otra persona.
\\
\\
La escucha activa es la práctica de prepararse para escuchar, observar los mensajes verbales y no verbales que se están enviando, y luego proporcionar retroalimentación apropiada con el fin de mostrar atención al mensaje que se está presentando \cite{UNIR2024}.

\item \textbf{Gamificación:} Aplicación de estrategias y mecánicas de juegos en contextos no jugables, con el fin de identificar aspectos susceptibles de ser convertidos en juegos o dinámicas lúdicas dentro de una actividad, tarea o mensaje determinado en un entorno de NO-juego. El objetivo es crear una experiencia significativa y motivadora que incentive un cambio de comportamiento, transmita un mensaje o contenido, y promueva una vinculación especial con los usuarios. En resumen, gamificar implica aplicar pensamientos y técnicas de juego para que las personas adopten ciertos comportamientos en contextos no relacionados con los juegos \cite{jenui2014}.

\item \textbf{Experiencia de usuario (UX):} La Experiencia de Usuario (UX) se centra en la percepción y sensaciones que experimenta una persona al interactuar con un producto, servicio o plataforma digital. Esta experiencia engloba todos los elementos de interacción, desde la forma en que se navega y se visualiza, hasta la facilidad con la que se utiliza y la satisfacción que genera en el usuario. En resumen, se trata de la impresión general y el nivel de comodidad que un usuario experimenta al utilizar un sitio web, una aplicación móvil u otro producto digital. Además, la UX no solo se limita a la utilidad y la estética, sino que también se preocupa por la emoción y el impacto emocional que deja en el usuario durante y después de la interacción.

\item \textbf{Plataformas de e-learning:} Una plataforma de e-learning es un entorno virtual diseñado para facilitar la administración y el acceso a una variedad de recursos educativos en línea. Estos recursos incluyen cursos completos, materiales multimedia, herramientas de comunicación interactivas, evaluaciones y la capacidad de monitorear el progreso del estudiante a lo largo del tiempo. Además de proporcionar contenido educativo, estas plataformas suelen ofrecer funciones adicionales como foros de discusión, chats en tiempo real, y la posibilidad de realizar trabajos colaborativos. 
\\
\\
Las plataformas de e-learning brindan una experiencia integral de aprendizaje en línea, permitiendo a los estudiantes aprender de manera flexible y a los educadores gestionar eficazmente el proceso educativo.

\item \textbf{Retroalimentación:} La retroalimentación es un proceso mediante el cual se proporciona información sobre el desempeño de una persona o sistema, con el propósito de mejorar o corregir su comportamiento o resultados. Esta información se entrega después de haber realizado una acción o tarea, y permite evaluar el desempeño para realizar ajustes que conduzcan a una mejora futura. La retroalimentación puede ser tanto positiva, al reconocer y elogiar un buen desempeño, como negativa, al señalar errores o deficiencias. Su objetivo es proporcionar información útil para el desarrollo y crecimiento de la persona o sistema que la recibe. Es una herramienta valiosa en diversos contextos, como el educativo, laboral o personal, ya que permite identificar áreas de mejora y tomar medidas para mejorar el desempeño en el futuro \cite{Escuela2024}.

\item \textbf{Roadmap de aprendizaje:} Un roadmap es una representación visual de alto nivel que muestra los objetivos que se desean alcanzar y las etapas necesarias para lograrlos, permitiendo planificar, priorizar y comunicar de manera clara la dirección de un proyecto \cite{MiroRoadmap2025}. En el contexto del aprendizaje, un roadmap se adapta como una hoja de ruta formativa que organiza metas educativas, actividades, recursos y plazos, facilitando que los estudiantes visualicen su progreso, estructuren sus procesos de estudio y tomen decisiones informadas sobre cómo avanzar de manera autónoma en su desarrollo.

\item \textbf{Meta de aprendizaje:} Una meta de aprendizaje representa un resultado deseado a largo plazo, normalmente asociado al desarrollo de una habilidad o competencia. Las metas funcionan como un marco de referencia que guía la planificación del estudio, la toma de decisiones y la motivación sostenida durante el proceso educativo.

\item \textbf{Objetivo:} Un objetivo de aprendizaje describe un resultado específico y alcanzable que el estudiante pretende lograr dentro de un periodo determinado. A diferencia de las metas (que son más amplias y de largo plazo), los objetivos desglosan propósitos generales en avances concretos, medibles y orientados a la acción.

\item \textbf{Tarea:} Una tarea es una acción puntual y operativa que el estudiante debe realizar para avanzar hacia el cumplimiento de un objetivo. Las tareas representan el nivel más específico dentro de la planificación del aprendizaje, ya que indican actividades concretas que permiten materializar el progreso.

\item \textbf{Agente de inteligencia artificial:} Sistema computacional capaz de percibir su entorno, procesar la información disponible y tomar decisiones de manera autónoma con el fin de alcanzar uno o varios objetivos definidos.

\item \textbf{Prompting:} Es el proceso de diseñar, estructurar y optimizar las instrucciones proporcionadas a un modelo de inteligencia artificial generativa (especialmente a modelos de lenguaje de gran escala) con el fin de dirigir su output y mejorar la calidad de las respuestas.

\item \textbf{LangChain:} Es un marco de trabajo de código abierto diseñado para facilitar la construcción de aplicaciones que integran modelos de lenguaje de gran escala con diversas fuentes de información externas. Su propósito es permitir que estos modelos no solo generen texto, si no que interactúen de manera estructurada con base de datos, documentos, APIs o sistemas internos, ampliando sus capacidades más allá del conocimiento implícito con el que fueron entrenados.

\item \textbf{LangGraph:} Es un marco de trabajo de código abierto orientado al desarrollo y orquestación de flujos de trabajo basados en agentes de inteligencia artificial. Su arquitectura se fundamenta en la representación mediante grafos, donde cada nodo corresponde a una etapa específica del proceso y las transiciones (aristas) definen la secuencia lógica o condicional entre dichas etapas. Este enfoque permite construir sistemas capaces de mantener estado, utilizar memoria contextual, repetir ciclos cuando es necesario y coordinar múltiples agentes que colaboran en la resolución de tareas complejas.

\item \textbf{Text-to-Speech (TTS):} Tecnología que transforma texto escrito en audio mediante algoritmos capaces de generar voz artificial de manera natural y comprensible.

\end{itemize}